% ============================================================================
% LC3 CC531 A — Presentación (10 minutos)
% Compilar: latexmk -pdf presentacion.tex
% Todas las cifras provienen de evidencias/salidas/ del repositorio.
% ============================================================================
\documentclass[10pt,aspectratio=169]{beamer}
\usepackage[utf8]{inputenc}
\usepackage[T1]{fontenc}
\usepackage[spanish]{babel}
\usepackage{booktabs}
\usepackage{tikz}
\usepackage{pgfplots}
\pgfplotsset{compat=1.18}

\usetheme{default}
\usecolortheme{beaver}
\setbeamertemplate{navigation symbols}{}
\setbeamertemplate{footline}{%
  \hfill{\scriptsize\color{gray}\insertframenumber\,/\,\inserttotalframenumber}\hspace*{6pt}\vspace*{4pt}}
\setbeamerfont{frametitle}{size=\large,series=\bfseries}

\definecolor{azul}{RGB}{20,60,110}
\definecolor{naranja}{RGB}{200,85,10}
\setbeamercolor{frametitle}{fg=azul}
\setbeamercolor{structure}{fg=azul}
\setbeamercolor{alerted text}{fg=naranja}
\setbeamertemplate{itemize items}[circle]

\newcommand{\fuente}[1]{{\scriptsize\color{gray}#1}}

\title[Desapariciones en el Perú 2019--2024]{Desapariciones de personas en el Perú (2019--2024)}
\subtitle{Dónde, quiénes y cuándo: 121,427 casos analizados con Spark y PostgreSQL}
\author[Pérez V. · Pezo]{Arbués Enrique Pérez Villanueva \and Sergio Pezo}
\institute[UNI]{CC531 A --- Análisis en Macrodatos (2026-2) · Universidad Nacional de Ingeniería}
\date{Laboratorio Calificado 03 · Noviembre 2026}

\begin{document}

% ---------------------------------------------------------------- 1 portada
\begin{frame}
  \titlepage
  \begin{center}
    \footnotesize
    \textit{Dataset: Denuncias por desaparición de personas [MININTER], 2019--2024}\\
    \textit{Plataforma Nacional de Datos Abiertos del Perú}
  \end{center}
\end{frame}

% ---------------------------------------------------------------- 2 problema
\begin{frame}{Qué analizamos y por qué importa}
  \begin{columns}[T]
    \begin{column}{0.50\textwidth}
      \textbf{La fuente}\\[3pt]
      \small Registro de denuncias por desaparición que publica el Ministerio del Interior
      en la Plataforma Nacional de Datos Abiertos (rubro Seguridad y Justicia).
      \vspace{6pt}
      \begin{itemize}\small
        \item 93,831 registros · \alert{121,427 personas}
        \item 13 campos · enero 2019 -- diciembre 2024
        \item Grano: fecha $\times$ ubigeo $\times$ perfil de la víctima
        \item Medida: \texttt{CANTIDAD} (personas por registro)
      \end{itemize}
      \vspace{4pt}
      \footnotesize\color{gray}Un solo CSV de 8.1 MB: se procesa en segundos en local[*].
    \end{column}
    \begin{column}{0.46\textwidth}
      \centering
      \begin{tikzpicture}[scale=0.78, every node/.style={font=\scriptsize}]
        \node[draw, fill=blue!8, rounded corners, text width=4.3cm, align=left, inner sep=5pt] {
          \textbf{\color{azul}Fecha}\\ AÑO, MES\\[3pt]
          \textbf{\color{azul}Geografía}\\ UBIGEO, DPTO, PROV, DIST\\[3pt]
          \textbf{\color{azul}Perfil}\\ SEXO, RANGO\_EDAD, EST\_CIVIL,\\ NACIONALIDAD, NIVEL\_EDUC., OCUPACIÓN\\[3pt]
          \textbf{\color{naranja}Medida}\\ CANTIDAD (1--18)
        };
      \end{tikzpicture}
    \end{column}
  \end{columns}
\end{frame}

% ---------------------------------------------------------------- 3 lo construido
\begin{frame}{Qué se construyó}
  \begin{columns}[T]
    \begin{column}{0.47\textwidth}
      \textbf{Tres capas de cómputo}\\[3pt]
      \begin{enumerate}\small
        \item \textbf{RDD + mapReduce}: 5 consultas (3a, 3b, 3c$\times$3)
        \item \textbf{API de DataFrame}: 14 consultas (5a--5d)
        \item \textbf{Vistas temporales + SQL}: 6 consultas (5e, 5f)
      \end{enumerate}
      \vspace{3pt}
      \small 25 consultas con objetivos distintos, controladas contra un mapa de consultas.
    \end{column}
    \begin{column}{0.49\textwidth}
      \textbf{Más la base de datos}\\[3pt]
      \begin{itemize}\small
        \item PostgreSQL 16 · \alert{modelo estrella} de 9 tablas
        \item ETL en tres pasos: staging $\to$ dimensiones $\to$ hechos
        \item 93,831 hechos cargados y verificados
        \item 6 libros de Excel generados \alert{con Spark}
      \end{itemize}
    \end{column}
  \end{columns}
  \vspace{8pt}
  \begin{center}
    \footnotesize\color{azul}
    Spark 3.3.2 · Scala 2.12.15 · OpenJDK 11 · PostgreSQL 16.16 · JDBC 42.7.6 · spark-excel 0.18.7
  \end{center}
\end{frame}

% ---------------------------------------------------------------- 4 el dato clave
\begin{frame}{El dato que cambia todas las cuentas}
  \begin{columns}[T]
    \begin{column}{0.44\textwidth}
      \centering
      \begin{tabular}{lr}
        \toprule
        Filas del CSV & 93,831\\
        \alert{Personas} & \alert{121,427}\\
        Filas con 1 persona & 78,128 (83.3\%)\\
        Promedio por registro & 1.29\\
        Mediana & 1\\
        Máximo en un hecho & 18\\
        \bottomrule
      \end{tabular}
      \\[6pt]
      \fuente{Consultas 3-b y 5-d.1}
    \end{column}
    \begin{column}{0.52\textwidth}
      \small
      \textbf{Un registro no es una persona.}\\
      Contar filas subestima el fenómeno en 27,596 personas: toda consulta agrega con
      \texttt{SUM(CANTIDAD)}.
      \vspace{6pt}

      \textbf{La cola es donde está el riesgo.}\\
      1,400 hechos con 5 o más víctimas (1.5\% de los registros) concentran
      \alert{9,188 personas} (7.6\%): allí deben activarse protocolos de búsqueda grupal.
      \vspace{6pt}

      \textbf{Faltan datos, pero no se inventan.}\\
      Nivel educativo y ocupación vacíos en el 65\% de las filas $\to$ categoría
      \texttt{NO ESPECIFICA}, sin imputar.
    \end{column}
  \end{columns}
\end{frame}

% ---------------------------------------------------------------- 5 serie anual
\begin{frame}{La serie anual: 2020 rompe la tendencia y no se recupera}
  \centering
  \begin{tikzpicture}
    \begin{axis}[
      width=13.4cm, height=5.2cm, ybar, bar width=22pt,
      ymin=0, ymax=29000, ylabel={\scriptsize personas},
      scaled y ticks=false, ytick={0,10000,20000}, yticklabels={0,10\,000,20\,000},
      xtick={2019,2020,2021,2022,2023,2024}, xtick=data,
      nodes near coords,
      nodes near coords style={font=\scriptsize,/pgf/number format/1000 sep={}},
      yticklabel style={/pgf/number format/1000 sep={}},
      axis line style={draw=gray!60}, tick style={draw=gray!60}, enlarge x limits=0.10
    ]
    \addplot[fill=azul!70] coordinates {(2019,26331) (2020,18481) (2021,20141) (2022,18862) (2023,18192) (2024,19420)};
    \end{axis}
  \end{tikzpicture}
  \\[2pt]
  \begin{columns}[T]
    \begin{column}{0.48\textwidth}
      \small\centering
      2020 cae \alert{29.8\%} frente a 2019:\\
      de 26,331 a 18,481 personas.
    \end{column}
    \begin{column}{0.48\textwidth}
      \small\centering
      Ningún año posterior vuelve al nivel inicial:\\
      2024 queda \alert{26.2\%} por debajo de 2019.
    \end{column}
  \end{columns}
  \vspace{2pt}
  \fuente{Fuente: consulta 5-e.3 (vistas temporales con join triple). Valores en personas, no en filas.}
\end{frame}

% ---------------------------------------------------------------- 6 hallazgo 1
\begin{frame}{Hallazgo 1: cuatro de cada diez son adolescentes mujeres}
  \begin{columns}[T]
    \begin{column}{0.50\textwidth}
      \centering
      \begin{tabular}{llrr}
        \toprule
        Grupo etario & Sexo & Personas & \% total\\
        \midrule
        \alert{Adolescente} & \alert{F} & \alert{51,021} & \alert{42.02}\\
        Adolescente & M & 13,309 & 10.96\\
        Adulto & M & 12,389 & 10.20\\
        Adulto joven & F & 12,042 & 9.92\\
        Adulto joven & M & 10,125 & 8.34\\
        Adulto mayor & M & 6,631 & 5.46\\
        \bottomrule
      \end{tabular}
      \\[3pt]
      \fuente{Consulta 5-e.2 · 18 combinaciones · base 121,427 personas}
    \end{column}
    \begin{column}{0.46\textwidth}
      \small
      El fenómeno tiene un perfil de \alert{género y edad} muy definido: las adolescentes
      mujeres concentran el \alert{42\%} de todos los casos del país.
      \vspace{5pt}

      La brecha con los adolescentes hombres es de casi \alert{4 a 1} (51,021 vs 13,309).
      \vspace{5pt}

      Los registros de mujeres tienen desviación 1.05 frente a 0.60 de los hombres: desaparecen
      en \alert{grupos}, no solas.
      \vspace{3pt}
      \fuente{Consulta 5-d.3}
    \end{column}
  \end{columns}
\end{frame}

% ---------------------------------------------------------------- 7 hallazgo 2
\begin{frame}{Hallazgo 2: el riesgo relativo de los menores crece con la ruralidad}
  \begin{columns}[T]
    \begin{column}{0.62\textwidth}
      \centering
      \begin{tikzpicture}
        \begin{axis}[
          width=8.6cm, height=5.8cm, xbar, bar width=9pt,
          xmin=45, xmax=78, xlabel={\scriptsize \% de menores de 18},
          symbolic y coords={LIMA MET.,AREQUIPA,PUNO,CALLAO,HUANCAV.,CUSCO,LA LIB.,JUN\'IN,PIURA,LAMBAYEQUE,HU\'ANUCO,CAJAMARCA,LORETO,UCAYALI,MADRE DE DIOS},
          ytick=data, y tick label style={font=\tiny},
          nodes near coords, nodes near coords style={font=\tiny},
          axis line style={draw=gray!60}, tick style={draw=gray!60}, enlarge y limits=0.04
        ]
        \addplot[fill=azul!60] coordinates {
         (50.02,LIMA MET.) (52.27,AREQUIPA) (52.60,PUNO) (54.81,CALLAO) (57.56,HUANCAV.)
         (59.59,CUSCO) (61.20,LA LIB.) (62.36,JUN\'IN) (62.38,PIURA) (65.90,LAMBAYEQUE)
         (68.50,HU\'ANUCO) (68.97,CAJAMARCA) (70.25,LORETO) (70.74,UCAYALI) (72.21,MADRE DE DIOS)};
        \end{axis}
      \end{tikzpicture}
      \\[1pt]
      \fuente{Consulta 3-c.1 · 26 departamentos · se muestran 15 para lectura}
    \end{column}
    \begin{column}{0.34\textwidth}
      \small
      \textbf{Lectura}\\[3pt]
      Madre de Dios: \alert{72.2\%} de las desapariciones son de menores (473 de 655).
      \vspace{5pt}

      Lima Metropolitana: \alert{50.6\%}.
      \vspace{5pt}

      Brecha: \alert{22.2 puntos porcentuales}.
      \vspace{8pt}

      Consecuencia: priorizar por \emph{riesgo relativo}, no solo por volumen absoluto.
    \end{column}
  \end{columns}
\end{frame}

% ---------------------------------------------------------------- 8 hallazgo 3
\begin{frame}{Hallazgo 3: diciembre y febrero concentran la demanda}
  \begin{columns}[T]
    \begin{column}{0.60\textwidth}
      \centering
      \begin{tikzpicture}
        \begin{axis}[
          width=8.8cm, height=5.4cm,
          ymin=6.8, ymax=10.1, ylabel={\scriptsize \% del total},
          xtick={1,2,3,4,5,6,7,8,9,10,11,12},
          xticklabels={E,F,M,A,M,J,J,A,S,O,N,D}, xtick=data,
          nodes near coords, nodes near coords style={font=\tiny},
          axis line style={draw=gray!60}, tick style={draw=gray!60}, enlarge x limits=0.05
        ]
        \addplot[mark=*, thick, color=naranja] coordinates {(1,7.89) (2,9.16) (3,8.81) (4,7.25) (5,7.72) (6,7.33) (7,8.17) (8,8.39) (9,8.39) (10,8.55) (11,8.85) (12,9.49)};
        \addplot[dashed, gray, domain=1:12] {8.333};
        \end{axis}
      \end{tikzpicture}
      \\[1pt]
      \fuente{Consulta 3-c.3 · línea punteada: 8.33\% esperado sin estacionalidad}
    \end{column}
    \begin{column}{0.36\textwidth}
      \small
      \textbf{Diciembre} 9.49\% (índice \alert{1.14}) y \textbf{febrero} 9.16\% (1.10) superan el
      esperado; \textbf{abril} 7.25\% (0.87) y \textbf{junio} 7.33\% (0.88) quedan por debajo.
      \vspace{5pt}

      Amplitud: \alert{2.24 puntos} entre el mes más alto y el más bajo.
      \vspace{5pt}

      Por trimestre: T4 26.9\%, T1 25.9\%, T3 25.0\%, T2 22.3\%.
    \end{column}
  \end{columns}
\end{frame}

% ---------------------------------------------------------------- 9 hallazgo 4
\begin{frame}{Hallazgo 4: la geografía del problema es estable}
  \begin{columns}[T]
    \begin{column}{0.52\textwidth}
      \centering
      \begin{tabular}{lrr}
        \toprule
        Distrito & Personas & \% país\\
        \midrule
        San Juan de Lurigancho & 4,775 & 3.9\\
        Lima & 3,153 & 2.6\\
        San Martín de Porres & 2,646 & 2.2\\
        Callao & 2,642 & 2.2\\
        Ate & 2,637 & 2.2\\
        Comas & 2,338 & 1.9\\
        Villa María del Triunfo & 2,053 & 1.7\\
        \bottomrule
      \end{tabular}
      \\[3pt]
      \fuente{Consultas 5-b.2 (2019--2024) y 5-f.1 (solo 2024)}
    \end{column}
    \begin{column}{0.44\textwidth}
      \small
      \textbf{Por región natural}\\[3pt]
      Costa: \alert{65.6\%} de las personas en 567 ubigeos.\\
      Sierra: 28.2\% pero en \alert{804 ubigeos}: más repartido.\\
      Selva: 6.1\% y la menor densidad por distrito.
      \vspace{6pt}

      El ranking de 2024 repite el acumulado, con Ate subiendo al 3.\textsuperscript{er} lugar:
      la focalización no cambia de un año a otro.
      \vspace{3pt}
      \fuente{Consultas 5-c.1 y 5-f.1}
    \end{column}
  \end{columns}
\end{frame}

% ---------------------------------------------------------------- 10 pandemia
\begin{frame}{2020: el trimestre que cortó la serie}
  \begin{columns}[T]
    \begin{column}{0.52\textwidth}
      \centering
      \begin{tabular}{lrr}
        \toprule
        Trimestre & Personas & vs. anterior\\
        \midrule
        2019-T3 & 6,628 & --\\
        2019-T4 & \alert{7,164} & +8.1\%\\
        2020-T1 & 6,018 & $-16.0$\%\\
        2020-T2 & \alert{2,485} & \alert{$-58.7$\%}\\
        2020-T3 & 4,564 & +83.7\%\\
        2024-T4 & 5,613 & --\\
        \bottomrule
      \end{tabular}
      \\[3pt]
      \fuente{Consulta 5-d.4 · fecha construida con to\_date + quarter}
    \end{column}
    \begin{column}{0.44\textwidth}
      \small
      El segundo trimestre de 2020 cae \alert{65\%} respecto de 2019-T4 (7,164 $\to$ 2,485).
      \vspace{6pt}

      La recuperación es parcial: el mejor trimestre reciente (2024-T4, 5,613) sigue
      \alert{22\%} por debajo de 2019-T4.
      \vspace{6pt}

      De los 60 meses comparables contra su año anterior, 21 suben y 39 bajan: caídas
      frecuentes, repuntes más fuertes.
      \vspace{3pt}
      \fuente{Consulta 5-e.1}
    \end{column}
  \end{columns}
\end{frame}

% ---------------------------------------------------------------- 11 ETL
\begin{frame}{Modelo estrella y ETL en PostgreSQL 16}
  \begin{columns}[T]
    \begin{column}{0.42\textwidth}
      \centering
      \begin{tikzpicture}[scale=0.60, every node/.style={font=\tiny},
        dim/.style={draw, fill=blue!8, rounded corners=1pt, inner sep=1.5pt, text width=2.1cm, align=center},
        fac/.style={draw, fill=orange!15, rounded corners=1pt, inner sep=2.5pt, text width=2.5cm, align=center}]
        \node[fac] (h) at (0,0) {\textbf{hecho\_desaparicion}\\ 93,831 filas\\ cantidad};
        \node[dim] (t) at (-3.1,1.6) {dim\_tiempo};
        \node[dim] (g) at (3.1,1.6) {dim\_geografia\\ 1,506};
        \node[dim] (e) at (-3.1,0.55) {dim\_rango\_edad};
        \node[dim] (s) at (-3.1,-0.5) {dim\_sexo};
        \node[dim] (c) at (3.1,0.55) {dim\_estado\_civil};
        \node[dim] (n) at (3.1,-0.5) {dim\_nacionalidad};
        \node[dim] (v) at (-3.1,-1.55) {dim\_nivel\_educ.};
        \node[dim] (o) at (3.1,-1.55) {dim\_ocupacion};
        \draw[-latex, gray] (t) -- (h); \draw[-latex, gray] (g) -- (h);
        \draw[-latex, gray] (e) -- (h); \draw[-latex, gray] (s) -- (h);
        \draw[-latex, gray] (c) -- (h); \draw[-latex, gray] (n) -- (h);
        \draw[-latex, gray] (v) -- (h); \draw[-latex, gray] (o) -- (h);
      \end{tikzpicture}
    \end{column}
    \begin{column}{0.55\textwidth}
      \small
      \textbf{1. Extracción}: \texttt{\textbackslash copy} del CSV a staging (93,831 filas).\\
      \vspace{3pt}
      \textbf{2. Transformación}: vacíos $\to$ \texttt{NO ESPECIFICA}; ubigeo a 6 dígitos;
      región natural, periodo y grupo etario derivados.\\
      \vspace{3pt}
      \textbf{3. Llenado}: dimensiones por \texttt{DISTINCT}, hechos con \texttt{cantidad > 0},
      todo en una transacción.\\
      \vspace{6pt}
      \alert{Calidad}: el ubigeo 150121 aparece con dos nombres (PUEBLO LIBRE y su antiguo
      MAGDALENA VIEJA); sin deduplicar, los joins sumaban 82 personas de más.
      \vspace{3pt}
      \fuente{SQL en sql/01\_ddl.sql y sql/02\_etl.sql}
    \end{column}
  \end{columns}
\end{frame}

% ---------------------------------------------------------------- 12 validacion
\begin{frame}{Validación: la misma cifra en tres tecnologías}
  \begin{center}
    \begin{tabular}{llr}
      \toprule
      Capa & Consulta de control & Personas\\
      \midrule
      RDD / mapReduce & 3-b (suma y estadísticos) & \alert{121,427}\\
      DataFrame API & 5-c.1 (join con dimensión geográfica) & \alert{121,427}\\
      PostgreSQL 16 & \texttt{SUM(cantidad)} tras el ETL & \alert{121,427}\\
      Dump restaurado & base nueva desde \texttt{dump-denuncias.sql} & \alert{121,427}\\
      \bottomrule
    \end{tabular}
  \end{center}
  \vspace{8pt}
  \begin{columns}[T]
    \begin{column}{0.48\textwidth}
      \small
      \textbf{Mismo orden en Spark y en SQL:}\\
      Lima Metropolitana 39,490 · Junín 7,256 · Cusco 7,214 · Arequipa 6,633 ·
      Lambayeque 6,231.
    \end{column}
    \begin{column}{0.48\textwidth}
      \small
      \textbf{Joins externos como control:}\\
      LEFT: 0 ubigeos sin registros. RIGHT: 0 meses sin datos. FULL OUTER: 26 departamentos con
      víctimas menores y adultas.
    \end{column}
  \end{columns}
\end{frame}

% ---------------------------------------------------------------- 13 conclusiones
\begin{frame}{Conclusiones}
  \begin{enumerate}\small
    \item \textbf{Perfil definido}: adolescentes mujeres, \alert{42\%} del total nacional, con
          hechos de múltiples víctimas más frecuentes que en hombres.
    \item \textbf{Riesgo relativo y ruralidad}: 72.2\% en Madre de Dios frente a 50.6\% en Lima;
          la priorización debería usar esa razón, no solo el volumen.
    \item \textbf{2020 cortó la serie} y no se recuperó: caída de 65\% en un trimestre y 2024
          todavía 22\% por debajo de 2019-T4.
    \item \textbf{Geografía estable}: los mismos distritos de Lima Norte y Este encabezan el
          ranking en 2024 y en el acumulado.
    \item \textbf{Técnicamente}: las tres capas cuadran en 121,427 personas y el modelo estrella
          reproduce consulta por consulta los resultados de Spark.
  \end{enumerate}
\end{frame}

% ---------------------------------------------------------------- 14 referencias
\begin{frame}{Referencias y reproducibilidad}
  \small
  \begin{itemize}
    \item Ministerio del Interior del Perú. \textit{Denuncias por desaparición de personas},
          2019--2024. Plataforma Nacional de Datos Abiertos.\\[2pt]
          {\scriptsize\url{https://www.datosabiertos.gob.pe/dataset/denuncias-por-desaparici\%C3\%B3n-de-personas-mininter}}
    \item Zaharia, S. et al. \textit{Resilient Distributed Datasets: A Fault-Tolerant Abstraction
          for In-Memory Cluster Computing}. NSDI, 2012.
    \item Armbrust, M. et al. \textit{Spark SQL: Relational Data Processing in Spark}. SIGMOD, 2015.
    \item Kimball, R., Ross, M. \textit{The Data Warehouse Toolkit}. 3.\textsuperscript{a} ed., Wiley, 2013.
    \item Apache Software Foundation. \textit{Spark 3.3.2 Documentation} y
          \textit{PostgreSQL 16 Documentation}.
  \end{itemize}
  \vspace{8pt}
  \begin{center}
    \footnotesize
    Repositorio: \url{https://github.com/Arbues/lc3-cc531-desapariciones}\\
    \texttt{bin/ejecutar.sh <consulta.scala>} · \texttt{bin/correr\_todas.sh} regenera las 26 salidas
  \end{center}
\end{frame}

% ---------------------------------------------------------------- 15 cierre
\begin{frame}{Gracias}
  \centering
  \vspace{20pt}
  {\Large Preguntas}\\[14pt]
  \small
  Arbués Enrique Pérez Villanueva · Sergio Pezo\\
  CC531 A --- Análisis en Macrodatos · Facultad de Ciencias, UNI\\[10pt]
  \textit{121,427 personas, 1,506 distritos, 6 años: el patrón se repite}
\end{frame}

\end{document}
